\begin{abstract}
\noindent
Grokking is the phenomenon in which a small neural network trained on a
synthetic algorithmic task first memorises the training set and then,
after thousands of additional training steps with no apparent progress,
abruptly generalises. \citet{nanda2023progress} showed that during the
apparent plateau a one-layer transformer is silently constructing a
discrete-Fourier-transform algorithm for modular addition. We extend
that analysis to a \emph{multi-task} setting in which a single one-layer
transformer is trained jointly on three modular arithmetic tasks ---
addition, subtraction, and multiplication modulo a prime $p$. Two of
these tasks live on the additive group $\Z/p\Z$ and one lives on the
multiplicative group $(\Z/p\Z)^*$, which has different Fourier-transform
structure.

We replicate the single-task addition baseline of
\citet{nanda2023progress} ($p = 113$, $\mathrm{train\_frac} = 0.30$,
AdamW with $\mathrm{weight\_decay} = 1.0$, grokking at step
$\sim$$7{,}200$) and obtain analogous sharp grokking transitions for
single-task subtraction and multiplication. Joint training on
\textsc{add}+\textsc{sub} grokes \emph{faster} than either single-task
baseline, suggesting that the shared additive Fourier basis is
discovered more quickly under joint gradient pressure. Joint training on
\textsc{add}+\textsc{sub}+\textsc{mul} produces a qualitatively
different dynamic: instead of a sharp plateau-then-collapse transition,
test accuracy on each task rises smoothly through 75{,}000 training
steps, reaching $0.60 / 0.48 / 0.32$ on the three tasks but never
crossing the 0.95 grok threshold. We argue this slow-ramp dynamic
reflects a residual-stream capacity competition between the two
distinct Fourier circuits that the additive and multiplicative tasks
each require. Code, configs, checkpoints, and all per-step metric
histories are released for reproducibility and continued analysis.
\end{abstract}
